A wave is driven by a random force that remembers its past. How fast can
its second moment grow? The answer depends on the covariance and on what the
stochastic product means. This chapter introduces the Stratonovich model and
states the answer in two parts: the scale of the growth, and the constant
at that scale. Both parts hold in space dimensions one, two and three.

We begin with a guitar in a sandstorm (Figure~\ref{Fig:walsh-guitar})
and a short history. Section~\ref{S:model} states the model, the
covariance and the initial data. Section~\ref{S:known} records what was
known, including the difference between the author and arXiv versions
of the source. Section~\ref{S:targets} states the two main theorems, and
Section~\ref{S:why-kappa} explains the scale. Figure~\ref{Fig:chapter-roadmap}
maps the rest of the book.

No earlier chapter is needed. On a first reading, the literature
comparison in Section~\ref{S:literature-context} and the outlook
in Sections~\ref{S:future-work} and~\ref{S:theory-outlook} can
wait. A reader entering later should start with the model and
initial data in Section~\ref{S:model}, Hypothesis~\ref{H:homog} and
the exponent in~\eqref{E:kappa}. Episodes~\VideoRef{1}{wx57PQ1iMJE}
and~\VideoRef{2}{MvNCNOq9rDw} of the video series follow this chapter. The
series retells the book in eleven short episodes, and each chapter guide
names its episode.
